% pipeline_overview.tex
% Two-page academic overview of the audio restoration pipeline.
% Compile with: pdflatex pipeline_overview.tex  (single pass, no bibtex needed)
\documentclass[9pt,twocolumn]{article}

% ── Geometry ─────────────────────────────────────────────────────────────────
\usepackage[top=0.60in,bottom=0.60in,left=0.55in,right=0.55in,
            columnsep=0.24in]{geometry}

% ── Fonts & encoding ─────────────────────────────────────────────────────────
\usepackage[T1]{fontenc}
\usepackage[utf8]{inputenc}
\usepackage{times}
\usepackage{mathptmx}
\usepackage{amssymb}
\usepackage{microtype}

% ── Tables ───────────────────────────────────────────────────────────────────
\usepackage{booktabs}
\usepackage{tabularx}

% ── Lists ────────────────────────────────────────────────────────────────────
\usepackage{enumitem}
\setlist{noitemsep,topsep=2pt,parsep=0pt,partopsep=0pt}

% ── Colors — [table] loads colortbl so \cellcolor works ──────────────────────
\usepackage[table]{xcolor}
\definecolor{cData}{RGB}{52,110,180}
\definecolor{cDist}{RGB}{198,70,55}
\definecolor{cDet} {RGB}{45,148,90}
\definecolor{cInp} {RGB}{130,75,185}
\definecolor{cEval}{RGB}{195,140,35}
\definecolor{cCut} {RGB}{210,60,50}

% ── TikZ ─────────────────────────────────────────────────────────────────────
\usepackage{tikz}
\usetikzlibrary{arrows.meta,calc}

% ── Section spacing ──────────────────────────────────────────────────────────
\usepackage{titlesec}
\titlespacing*{\section}{0pt}{5pt}{2pt}
\titleformat{\section}{\normalfont\bfseries\small}{\thesection}{0.4em}{}

% ── Paragraph spacing ────────────────────────────────────────────────────────
\setlength{\parskip}{3pt}
\setlength{\parindent}{7pt}

% ── References ───────────────────────────────────────────────────────────────
\usepackage[numbers,compress]{natbib}
\setlength{\bibsep}{2pt}

% ── Hyperlinks ───────────────────────────────────────────────────────────────
\usepackage[hidelinks]{hyperref}

% ── Allow slight line stretching to eliminate overfull hboxes ────────────────
\setlength{\emergencystretch}{3em}

% =============================================================================
\begin{document}
\pagestyle{empty}

% ── Full-width header: title + abstract + diagram ────────────────────────────
\twocolumn[{%
\begin{center}
{\large\bfseries
  Blind Detection and Acoustic Inpainting for Telephony Speech Restoration}\\[2pt]
{\small MIT MBAn GenAI Lab\quad|\quad April 2026}
\end{center}
\vspace{2pt}
\noindent\textbf{Abstract.}\;{\small
We present an end-to-end pipeline for restoring telephony speech corrupted by
packet-loss gaps without prior knowledge of cut locations.
Using the Harper Valley banking corpus~\cite{harpervalley}, SHA256-seeded
cuts are placed exclusively in VAD-detected speech regions under two fill
modes: hard zeros (baseline) and RFC\,3389 comfort-noise (realistic).
Two blind detectors---a sample-level silence scanner and a
half-wave-rectified spectral-flux detector~\cite{bello2005}---recover
timestamps.  Gaps are repaired via Qwen3-ASR~\cite{qwen3}\,/\,%
Whisper~\cite{whisper}, Qwen3-32B text correction, and
PlayDiffusion~\cite{playdiffusion} with duration-aware routing.
On 125 synthetic samples, zero-fill Tiers A--D achieve P\,=\,R\,=\,1.00;
comfort-noise Tier~E (P\,=\,0.03) exposes the signal-processing ceiling.
Word-error rate (WER) complements correct-action rate (CAR) in evaluation.
}
\vspace{3pt}

% ── Inline pipeline diagram ──────────────────────────────────────────────────
\begin{center}
\vspace{-2pt}
\begin{tikzpicture}[
  every node/.style={font=\scriptsize\sffamily, align=center},
  box/.style ={rectangle, rounded corners=2pt, draw, thick,
               text centered, inner sep=2.5pt, minimum height=0.46cm},
  arr/.style ={-{Stealth[length=3.8pt,width=3pt]}, thick},
  darr/.style={-{Stealth[length=3.8pt,width=3pt]}, thick, dashed, opacity=0.6},
  x=1cm, y=1cm,
]

%% Row 1 (y=0) — data + distortion (symmetric around x=0)
\node[box, minimum width=2.45cm, fill=cData!18, draw=cData!70]
  (hv)  at (-3.60,  0.00) {Harper Valley\\\tiny 8 tasks · 16\,kHz · mono};
\node[box, minimum width=2.80cm, fill=cDist!15, draw=cDist!60]
  (vad) at ( 0.00,  0.00) {VAD-Aware Distortion\\\tiny 20\,ms · adaptive\,$\tau$};
\node[box, minimum width=2.45cm, fill=cDist!25, draw=cDist!70]
  (da)  at ( 3.60,  0.00) {Distorted Audio\\\tiny 4 cuts · 1–200\,ms · SHA256};

\draw[arr, cDist!70] (hv)  -- (vad);
\draw[arr, cDist!70] (vad) -- (da);

%% Row 2 (y=-1.10) — blind detectors (symmetric ±1.8)
\node[box, minimum width=2.45cm, fill=cDet!18, draw=cDet!65]
  (sil) at (-1.80, -1.10) {Silence Detector\\\tiny $|x[n]|{<}10^{-4}$\,·\,pre-ctx};
\node[box, minimum width=2.45cm, fill=cDet!18, draw=cDet!65]
  (sf)  at ( 1.80, -1.10) {Spectral-Flux\\\tiny HWR-STFT\,·\,99.5th\,pct};

\coordinate (dafork) at (0.00, -0.57);
\draw[arr, cDet!65] (da.south) |- (dafork);
\draw[arr, cDet!65] (dafork)   -| (sil.north);
\draw[arr, cDet!65] (dafork)   -| (sf.north);

%% Row 3 (y=-2.18) — JSON timestamps
\node[box, minimum width=2.80cm, fill=cDet!30, draw=cDet!75]
  (json) at (0.00, -2.18) {Cut Timestamps (JSON)\\\tiny start\_sec · end\_sec · type};

\draw[arr, cDet!75] (sil.south) |- (json.west);
\draw[arr, cDet!75] (sf.south)  |- (json.east);

%% Row 4 (y=-3.26) — ASR + LLM (symmetric ±1.8)
\node[box, minimum width=2.45cm, fill=cInp!15, draw=cInp!60]
  (asr) at (-1.80, -3.26) {Qwen3-ASR-1.7B\\\tiny words + timestamps};
\node[box, minimum width=2.45cm, fill=cInp!15, draw=cInp!60]
  (llm) at ( 1.80, -3.26) {Qwen3-32B (LLM)\\\tiny transcript correction};

\coordinate (jsonfork) at (0.00, -2.73);
\draw[arr, cInp!60] (json.south) -- (jsonfork);
\draw[arr, cInp!60] (jsonfork)   -| (asr.north);
\draw[arr, cInp!60] (jsonfork)   -| (llm.north);

%% Row 5 (y=-4.34) — inpainting
\node[box, minimum width=5.10cm, fill=cInp!28, draw=cInp!70]
  (pd) at (0.00, -4.34)
  {PlayDiffusion Inpainting\\\tiny
   silence$\!\to\!$crossfade\;$|$\;
   glitch$\!\to\!$text-diff\;$|$\;
   normal$\!\to\!$InpaintInput};

\draw[arr, cInp!70] (asr.south) |- (pd.west);
\draw[arr, cInp!70] (llm.south) |- (pd.east);

%% Row 6 (y=-5.40) — evaluation
\node[box, minimum width=5.10cm, fill=cEval!22, draw=cEval!70]
  (eval) at (0.00, -5.40)
  {Agent Evaluation\\\tiny
   \textcolor{cData!80}{\textbf{clean}}\;/\;
   \textcolor{cDist!80}{\textbf{distorted}}\;/\;
   \textcolor{cInp!80}{\textbf{restored}}\quad$\to$\quad
   correct-action rate (CAR)};

\draw[arr, cEval!70] (pd.south) -- (eval.north);

%% Bypass arrows — run along outer edges, clear of all nodes
\draw[darr, cData!70]
  (hv.south) -- ++(0,-5.17) -| (eval.west);
\draw[darr, cDist!70]
  (da.south) -- ++(0,-5.17) -| (eval.east);

%% Legend
\node[font=\tiny\sffamily, anchor=north west, text=gray!75]
  at (-3.80, -5.88)
  {\textcolor{cData!80}{$\blacksquare$}\,data\enspace
   \textcolor{cDist!80}{$\blacksquare$}\,distortion\enspace
   \textcolor{cDet!80}{$\blacksquare$}\,detection\enspace
   \textcolor{cInp!80}{$\blacksquare$}\,inpainting\enspace
   \textcolor{cEval!80}{$\blacksquare$}\,evaluation\enspace
   \textcolor{gray!60}{$\cdots$}\,bypass};

\end{tikzpicture}
\end{center}
{\scriptsize\textbf{Figure\,1.}
End-to-end pipeline.
Dashed arrows carry
\textcolor{cData!80}{\textbf{clean}} and
\textcolor{cDist!80}{\textbf{distorted}} audio directly to evaluation
for three-condition comparison.}
\vspace{4pt}
}]% end \twocolumn header

% =============================================================================
% Two-column body
% =============================================================================

\section{Dataset and Distortion Model}
\label{sec:distortion}

The \textbf{Gridspace Stanford Harper Valley} corpus~\cite{harpervalley}
provides real telephony conversations for 8 banking actions at 16\,kHz mono.
\texttt{generate\_distorted\_audio.py} first runs an energy-based VAD on 20\,ms
frames, labelling speech whenever RMS exceeds
$\tau{=}\hat\mu_\mathrm{ns}+0.02(\hat\sigma_\mathrm{pk}-\hat\mu_\mathrm{ns})$
(20th-percentile noise floor plus 2\% of the dynamic range).
Four cuts per file, drawn uniformly from $[1,200]$\,ms, are placed
\emph{exclusively inside speech regions}: Harper Valley audio is
${\approx}$10\% speech, so blind placement would mostly hit silence.
Distortion is seeded by SHA256(\textit{stem})[:8], ensuring full
determinism without a manifest.
A second fill mode (\texttt{--fill-mode\,comfort-noise}) replaces zeroed
samples with bandpass-filtered Gaussian noise ($300$--$3\,400$\,Hz) scaled
to the local noise floor via a 20\,ms context window---modelling RFC\,3389
comfort-noise injection and cosine-faded at boundaries to prevent clicks.

% ── Figure 2 floats to the top of the next page ──────────────────────────────
\begin{figure*}[!t]
\centering
\begin{tikzpicture}[
  every node/.style={font=\scriptsize\sffamily, align=left},
  x=1cm, y=1cm,
]

%%──────────────────────────────────────────────────────────────────────────────
%% PANEL (i): Cut in inter-word silence (y-offset = 0)
%%──────────────────────────────────────────────────────────────────────────────
\begin{scope}[yshift=0cm]
  %% Time axis
  \draw[->, gray!45, thin] (-0.15, 0) -- (10.3, 0)
    node[right, font=\tiny\sffamily, text=gray!60] {$t$};

  %% Speech "balance" (x = 0.2 → 3.0)
  \draw[cData!70, line width=0.8pt]
    plot[smooth, tension=0.75] coordinates {
      (0.20, 0.00) (0.50, 0.31) (0.85,-0.29) (1.20, 0.34)
      (1.55,-0.27) (1.90, 0.33) (2.25,-0.30) (2.60, 0.26)
      (2.85,-0.16) (3.00, 0.00)
    };

  %% Natural pause before cut (x = 3.0 → 3.5), near-flat
  \draw[gray!50, line width=0.6pt]
    plot[smooth] coordinates {
      (3.00, 0.00) (3.15, 0.02) (3.35,-0.02) (3.50, 0.00)
    };

  %% CUT region (x = 3.5 → 4.7): red filled box
  \fill[cCut!17] (3.50,-0.52) rectangle (4.70, 0.52);
  \draw[cCut!65, thick, dashed] (3.50,-0.52) rectangle (4.70, 0.52);

  %% Natural pause after cut (x = 4.7 → 5.2)
  \draw[gray!50, line width=0.6pt]
    plot[smooth] coordinates {
      (4.70, 0.00) (4.85, 0.02) (5.05,-0.02) (5.20, 0.00)
    };

  %% Speech "account" (x = 5.2 → 9.8)
  \draw[cData!70, line width=0.8pt]
    plot[smooth, tension=0.75] coordinates {
      (5.20, 0.00) (5.55,-0.30) (5.90, 0.33) (6.25,-0.28)
      (6.60, 0.34) (6.95,-0.27) (7.30, 0.35) (7.65,-0.26)
      (8.00, 0.31) (8.35,-0.24) (8.70, 0.27) (9.10,-0.18)
      (9.50, 0.11) (9.80, 0.00)
    };

  %% Word-boundary dashed lines
  \draw[gray!30, thin, dashed] (3.00,-0.56) -- (3.00, 0.64);
  \draw[gray!30, thin, dashed] (5.20,-0.56) -- (5.20, 0.64);

  %% Labels above
  \node[above, text=cData!80, font=\tiny\sffamily\itshape]
    at (1.60, 0.38) {``balance''};
  \node[above, text=cData!80, font=\tiny\sffamily\itshape]
    at (7.50, 0.38) {``account''};
  \node[above, text=cCut!80,  font=\tiny\sffamily\bfseries, align=center]
    at (4.10, 0.53) {cut};

  %% Annotations below
  \node[below, text=gray!65, font=\tiny\sffamily] at (4.10,-0.57)
    {inter-word silence (within VAD speech region)};
  \node[text=cDet!75, font=\tiny\sffamily, anchor=west]
    at (0.00,-0.95)
    {Detection: pre-context check may fire → event suppressed as natural pause};
  \node[text=cInp!75, font=\tiny\sffamily, anchor=west]
    at (0.00,-1.25)
    {Routing: \textbf{crossfade} (no ASR word overlap at gap)};

  %% Panel label
  \node[font=\small\sffamily\bfseries, text=black, anchor=east]
    at (-0.15, 0) {(i)};
\end{scope}

%%──────────────────────────────────────────────────────────────────────────────
%% PANEL (ii): Cut at word boundary (y-offset = -2.5)
%%──────────────────────────────────────────────────────────────────────────────
\begin{scope}[yshift=-2.5cm]
  \draw[->, gray!45, thin] (-0.15, 0) -- (10.3, 0)
    node[right, font=\tiny\sffamily, text=gray!60] {$t$};

  %% Speech "check" (x = 0.2 → 3.8)
  \draw[cData!70, line width=0.8pt]
    plot[smooth, tension=0.75] coordinates {
      (0.20, 0.00) (0.55, 0.32) (0.90,-0.29) (1.25, 0.35)
      (1.60,-0.28) (1.95, 0.33) (2.30,-0.30) (2.65, 0.34)
      (3.00,-0.26) (3.35, 0.24) (3.65,-0.14) (3.80, 0.00)
    };

  %% CUT region straddles word boundary (x = 3.8 → 5.4)
  \fill[cCut!17] (3.80,-0.52) rectangle (5.40, 0.52);
  \draw[cCut!65, thick, dashed] (3.80,-0.52) rectangle (5.40, 0.52);

  %% Word boundary (x = 4.6) shown as line through cut box
  \draw[gray!35, thin, dashed] (4.60,-0.56) -- (4.60, 0.64);

  %% Speech "balance" (x = 5.4 → 9.8)
  \draw[cData!70, line width=0.8pt]
    plot[smooth, tension=0.75] coordinates {
      (5.40, 0.00) (5.75,-0.30) (6.10, 0.34) (6.45,-0.28)
      (6.80, 0.33) (7.15,-0.28) (7.50, 0.35) (7.85,-0.26)
      (8.20, 0.31) (8.55,-0.23) (8.90, 0.27) (9.25,-0.17)
      (9.55, 0.10) (9.80, 0.00)
    };

  %% Labels above
  \node[above, text=cData!80, font=\tiny\sffamily\itshape]
    at (2.00, 0.38) {``check''};
  \node[above, text=cData!80, font=\tiny\sffamily\itshape]
    at (7.60, 0.38) {``balance''};
  \node[above, text=cCut!80, font=\tiny\sffamily\bfseries, align=center]
    at (4.60, 0.53) {cut};
  \node[above, text=gray!55, font=\tiny\sffamily]
    at (4.60, 0.60) {\tiny word boundary};

  %% Annotations below
  \node[below, text=gray!65, font=\tiny\sffamily] at (4.60,-0.57)
    {cut straddles word boundary};
  \node[text=cDet!75, font=\tiny\sffamily, anchor=west]
    at (0.00,-0.95)
    {Detection: silence detector catches sudden amplitude drop};
  \node[text=cInp!75, font=\tiny\sffamily, anchor=west]
    at (0.00,-1.25)
    {Routing: \textbf{crossfade} (ASR word timestamps do not overlap gap centre)};

  \node[font=\small\sffamily\bfseries, text=black, anchor=east]
    at (-0.15, 0) {(ii)};
\end{scope}

%%──────────────────────────────────────────────────────────────────────────────
%% PANEL (iii): Cut mid-word (y-offset = -5.0)
%%──────────────────────────────────────────────────────────────────────────────
\begin{scope}[yshift=-5.0cm]
  \draw[->, gray!45, thin] (-0.15, 0) -- (10.3, 0)
    node[right, font=\tiny\sffamily, text=gray!60] {$t$};

  %% Speech "transfer" — first half (x = 0.2 → 2.8)
  \draw[cData!70, line width=0.8pt]
    plot[smooth, tension=0.75] coordinates {
      (0.20, 0.00) (0.55, 0.33) (0.90,-0.30) (1.25, 0.36)
      (1.60,-0.29) (1.95, 0.34) (2.30,-0.28) (2.60, 0.20)
      (2.80, 0.00)
    };

  %% CUT region mid-word (x = 2.8 → 4.2)
  \fill[cCut!17] (2.80,-0.52) rectangle (4.20, 0.52);
  \draw[cCut!65, thick, dashed] (2.80,-0.52) rectangle (4.20, 0.52);

  %% Speech "transfer" — second half (x = 4.2 → 6.4)
  \draw[cData!70, line width=0.8pt]
    plot[smooth, tension=0.75] coordinates {
      (4.20, 0.00) (4.55, 0.30) (4.90,-0.28) (5.25, 0.33)
      (5.60,-0.26) (5.95, 0.23) (6.20,-0.13) (6.40, 0.00)
    };

  %% Speech "money" (x = 7.2 → 9.8)
  \draw[cData!70, line width=0.8pt]
    plot[smooth, tension=0.75] coordinates {
      (7.20, 0.00) (7.55,-0.30) (7.90, 0.34) (8.25,-0.27)
      (8.60, 0.32) (8.95,-0.23) (9.30, 0.21) (9.60,-0.11)
      (9.80, 0.00)
    };

  %% Word-boundary markers for "transfer" token (full extent x=0.2 to x=6.4)
  \draw[gray!35, thin, dashed] (0.20,-0.56) -- (0.20, 0.76);
  \draw[gray!35, thin, dashed] (6.40,-0.56) -- (6.40, 0.76);

  %% Labels above
  %% "transfer" token label spans boundary markers; "cut" label above box
  \node[above, text=cData!80, font=\tiny\sffamily\itshape, align=center]
    at (3.30, 0.66) {``transfer'' (single ASR token)};
  \node[above, text=cCut!80, font=\tiny\sffamily\bfseries]
    at (3.50, 0.53) {cut};
  \node[above, text=cData!80, font=\tiny\sffamily\itshape]
    at (8.50, 0.38) {``money''};

  %% Annotation below
  \node[below, text=gray!65, font=\tiny\sffamily] at (3.50,-0.57)
    {cut lies entirely within one word token};
  \node[text=cDet!75, font=\tiny\sffamily, anchor=west]
    at (0.00,-0.95)
    {Detection: silence or spectral-flux detector (both applied)};
  \node[text=cInp!75, font=\tiny\sffamily, anchor=west]
    at (0.00,-1.25)
    {Routing (${\geq}$20\,ms): adaptive MFCC check $\to$ TTS regen \textbf{or} PlayDiffusion};

  \node[font=\small\sffamily\bfseries, text=black, anchor=east]
    at (-0.15, 0) {(iii)};
\end{scope}

\end{tikzpicture}
\vspace{4pt}
{\scriptsize\textbf{Figure\,2.}
The three cut location types relative to the ASR word lattice.
Blue curves are schematic speech waveforms; red dashed boxes mark the
removed (or zeroed) region; grey dashed verticals are word-token
boundaries.
Detection behaviour and inpainting routing differ for each type:
inter-word silence (i) may be suppressed by the pre-context check;
word-boundary cuts (ii) are caught but need only crossfade; mid-word
cuts (iii) require full neural inpainting when the gap exceeds 20\,ms.}
\end{figure*}

\section{Cut Taxonomy}
\label{sec:taxonomy}

Packet-loss cuts are not homogeneous.  They vary along two independent
dimensions---\emph{where} in the speech stream the gap falls and
\emph{how long} it lasts---and each combination requires a different
handling strategy.  A third axis, fill mode (\S\ref{sec:distortion}),
governs detectability.  Figure~2 illustrates the three location types;
Table~1 summarises routing across all combinations.

\noindent\textbf{Location.}
Cuts fall in one of three positions relative to ASR word timestamps.
\emph{(i) Inter-word silence:} the gap lies entirely within a natural
pause between adjacent words (still inside a VAD-active region).  No
phonetic content is removed.  The pre-context silence check may suppress
this event as a natural pause; if detected, crossfade suffices.
\emph{(ii) Word boundary:} the gap straddles the transition between two
words.  ASR places both adjacent tokens outside the gap, so no word
overlap is reported and crossfade is applied.
\emph{(iii) Mid-word:} the gap falls entirely within one word token.
ASR confirms overlap.  For gaps shorter than 20\,ms crossfade is used
unconditionally (perceptually transparent).  For longer gaps the adaptive
MFCC check (\S\ref{sec:inpainting}) selects between TTS regeneration
and PlayDiffusion inpainting.

\noindent\textbf{Duration.}
Cuts below 5\,ms (Tier\,C) are sub-threshold for most listeners; 5--20\,ms
(Tier\,B) produce a click while preserving intelligibility; 20--200\,ms
(Tier\,A) cause word-level loss.  Comfort-noise fills (Tier\,E, any
duration) defeat all signal-processing detectors because the fill is
spectrally indistinguishable from the telephony background.

\noindent\textbf{Table\,1 — Routing matrix.}
Crossfade (CF) handles the entire matrix except the single cell where
a mid-word gap exceeds 20\,ms and neural inpainting is needed.

\smallskip
{\footnotesize\setlength{\tabcolsep}{3pt}
\begin{tabularx}{\columnwidth}{lXXX}
\toprule
\textbf{Location} & \textbf{$<$5\,ms} & \textbf{5--20\,ms}
  & \textbf{$\geq$20\,ms} \\
\midrule
In silence
  & \cellcolor{cDet!14}CF$^{*}$
  & \cellcolor{cDet!14}CF$^{*}$
  & \cellcolor{cDet!14}CF$^{*}$ \\
Word boundary
  & \cellcolor{cDet!14}CF
  & \cellcolor{cDet!14}CF
  & \cellcolor{cDet!14}CF \\
Mid-word
  & \cellcolor{cDet!14}CF
  & \cellcolor{cDet!14}CF
  & \cellcolor{cInp!20}MFCC\,${\to}$\,TTS\,/\,PD \\
\bottomrule
\end{tabularx}
}
\vspace{1pt}
{\tiny CF\,=\,crossfade; PD\,=\,PlayDiffusion.\quad
$^{*}$Detector may suppress event as natural pause (pre-context check).}

\section{Blind Cut Detection}
\label{sec:detection}

\texttt{detect\_cuts.py} recovers timestamps without ground-truth access
using two complementary detectors.

\noindent\textbf{Silence detector.}
Contiguous samples with $|x[n]|\leq10^{-4}$ for $\geq0.5$\,ms are flagged.
A 5\,ms pre-context check discards regions where the preceding audio was
already quiet (natural pauses), retaining only sudden amplitude drops.

\noindent\textbf{Spectral-flux detector.}
STFT magnitude (10\,ms Hann window, 2\,ms hop) is computed frame-by-frame.
Half-wave-rectified flux
$\Phi_k{=}\sum_b\max(0,|M_{k,b}|{-}|M_{k{-}1,b}|)$
flags abrupt spectral increases~\cite{bello2005};
peaks above the 99.5th percentile within 10\,ms are merged.
Spectral peaks overlapping a silent region are absorbed.

\noindent\textbf{Evaluation} on 125 synthetic telephony-noise clips
(300--3\,400\,Hz, five tiers A--E by cut type):

{\footnotesize\setlength{\tabcolsep}{3pt}
\begin{tabularx}{\columnwidth}{lcccc}
\toprule
Tier & $n$ & Range & P & R / F\textsubscript{1} \\
\midrule
A easy   & 25 & 20--200\,ms & 1.000 & 1.000 \\
B medium & 25 &  5--20\,ms  & 1.000 & 1.000 \\
C hard   & 25 &   1--5\,ms  & 1.000 & 1.000 \\
D mixed  & 25 &  1--200\,ms & 1.000 & 1.000 \\
E cn     & 25 &  1--200\,ms & 0.029 & 0.103 \\
\midrule
\textbf{A--D} & 100 & --- & \textbf{1.000} & \textbf{1.000}\\
\bottomrule
\end{tabularx}}
\vspace{2pt}
{\tiny TP\,=\,257, FP\,=\,FN\,=\,0 (Tiers A--D, zero-fill).\;
Tier\,E (comfort-noise): fill is spectrally identical to telephony
background noise $\Rightarrow$ signal-processing detection ceiling;
neural methods required.}

\section{Audio Inpainting}
\label{sec:inpainting}

Each detected gap passes through three chained models (Google Colab, T4).

\noindent\textbf{Transcription.}
\textbf{Qwen3-ASR-1.7B}~\cite{qwen3} generates word-level timestamps via
forced alignment; \textbf{Whisper}~\cite{whisper} provides a fallback.

\noindent\textbf{Correction.}
The broken transcript is sent to \textbf{Qwen3-32B}~\cite{qwen3}
(OpenRouter) to infer the intended sentence near the cut.

\noindent\textbf{PlayDiffusion}~\cite{playdiffusion} fills gaps under
three routing rules (see also Table~1):
\begin{enumerate}[label=(\roman*)]
\item \emph{Short or silence gap:} duration ${<}20$\,ms \emph{or} no ASR
  word overlap; linear crossfade only---perceptually transparent and
  preserves speaker identity.
\item \emph{Acoustic glitch:} MFCC $L_2$ exceeds the per-file 95th-pct
  frame-distance baseline (adaptive threshold, not hardcoded); overlapping
  tokens are removed to force regeneration.
\item \emph{Normal mid-word gap:} \texttt{InpaintInput} with ASR-snapped
  word boundaries and LLM-corrected \texttt{output\_text}.
\end{enumerate}

\section{Agent Evaluation}
\label{sec:agent}

\texttt{eval\_runner.py} parses action labels from filenames
(\texttt{B03\_05.wav\,$\to$\,B03}) and scores 8 banking tasks across
three conditions.
\textbf{Correct-action rate (CAR):} Whisper~\cite{whisper} transcribes
each file; a keyword-scoring classifier assigns one of the 8 banking tasks.
Clean\,/\,distorted\,/\,restored audio set ceiling, floor, and system
under test: $\mathrm{CAR}(\mathrm{restored})>\mathrm{CAR}(\mathrm{distorted})$.
\textbf{WER} vs.\ clean references (stdlib edit-distance) traces the causal
chain: cuts$\to$WER$\uparrow\to$CAR$\downarrow$;
restoration$\to$WER$\downarrow\to$CAR$\uparrow$.

{\footnotesize\setlength{\bibsep}{2pt}
\begin{thebibliography}{9}
\bibitem{harpervalley}
E.~Lichtarge \textit{et al.}
Gridspace-Stanford Harper Valley Speech, 2020.
\bibitem{whisper}
A.~Radford \textit{et al.}
Robust speech recognition via large-scale weak supervision.
\textit{ICML}, 2023.
\bibitem{qwen3}
Qwen Team.
Qwen3 technical report. \textit{arXiv:2505.09388}, 2025.
\bibitem{playdiffusion}
PlayHT~Inc.
PlayDiffusion: neural audio inpainting, 2024.
\bibitem{bello2005}
J.\,P.~Bello \textit{et al.}
A tutorial on onset detection.
\textit{IEEE Trans.\ Speech Audio Process.}, 13(5):1035--1047, 2005.
\end{thebibliography}
}

\end{document}
